\textbf{Phần 1. Trắc nghiệm lựa chọn}
\setcounterpageref{ex}{0}

\begin{ex} %[3D1V59]
Số gồm $6$ chục nghìn, $4$ nghìn, $0$ trăm, $5$ chục và $2$ đơn vị được viết là
\choice
{$64\,502$}
{$64\,520$}
{\True $64\,052$}
{$60\,452$}
\loigiai{
Hàng chục nghìn là 6, nghìn là 4, trăm là 0, chục là 5, đơn vị là 2. Số đó là $64\,052$.
\textbf{Đáp án của câu hỏi là C}
}
\end{ex}

\begin{ex} %[3D1V59]
Số liền sau của số $99\,999$ là số nào?
\choice
{\True $100\,000$}
{$99\,998$}
{$100\,001$}
{$100\,010$}
\loigiai{
Số liền sau của 99 999 là: $99\,999 + 1 = 100\,000$.
\textbf{Đáp án của câu hỏi là A}
}
\end{ex}

\begin{ex} %[3D1V59]
Trong số $87\,354$, giá trị của chữ số 8 gấp giá trị của chữ số 7 bao nhiêu lần?
\choice
{Gấp 10 lần}
{Gấp 8 lần}
{Gấp 100 lần}
{\True Không chia hết, giá trị lần lượt là $80\,000$ và $7000$}
\loigiai{
Chữ số 8 ở hàng chục nghìn có giá trị $80\,000$. Chữ số 7 ở hàng nghìn có giá trị $7000$.
\textbf{Đáp án của câu hỏi là D}
}
\end{ex}

\begin{ex} %[3D1V59]
Từ năm chữ số $0; 2; 5; 7; 9$, viết được số có năm chữ số khác nhau bé nhất là
\choice
{$25\,079$}
{\True $20\,579$}
{$20\,759$}
{$20\,597$}
\loigiai{
Chữ số hàng chục nghìn bé nhất khác 0 là 2, tiếp theo là 0, 5, 7, 9. Số đó là $20\,579$.
\textbf{Đáp án của câu hỏi là B}
}
\end{ex}

\begin{ex} %[3D1V59]
Số chẵn lớn nhất có năm chữ số khác nhau là
\choice
{$99\,998$}
{$98\,765$}
{\True $98\,764$}
{$98\,762$}
\loigiai{
Các chữ số giảm dần từ hàng cao: 9, 8, 7, 6; chữ số hàng đơn vị chẵn khác nhau lớn nhất là 4. Số đó là $98\,764$.
\textbf{Đáp án của câu hỏi là C}
}
\end{ex}

\begin{ex} %[3D1V59]
Có bao nhiêu số tròn chục nghìn có năm chữ số bé hơn $100\,000$?
\choice
{\True $9$ số}
{$10$ số}
{$8$ số}
{$90$ số}
\loigiai{
Các số tròn chục nghìn: $10\,000, 20\,000, 30\,000, 40\,000, 50\,000, 60\,000, 70\,000, 80\,000, 90\,000$. Có tất cả 9 số.
\textbf{Đáp án của câu hỏi là A}
}
\end{ex}

\begin{ex} %[3D1V59]
Một số có năm chữ số có chữ số hàng chục nghìn bằng chữ số hàng đơn vị và bằng 5. Các chữ số còn lại đều bằng 0. Số đó là
\choice
{$50\,050$}
{$50\,500$}
{$55\,000$}
{\True $50\,005$}
\loigiai{
Hàng chục nghìn là 5, đơn vị là 5, các hàng khác là 0: $50\,005$.
\textbf{Đáp án của câu hỏi là D}
}
\end{ex}

\begin{ex} %[3D1V59]
Cho số $45\,678$. Nếu xóa đi chữ số 4 ở hàng chục nghìn thì số đó giảm đi bao nhiêu đơn vị?
\choice
{$4000$ đơn vị}
{\True $40\,000$ đơn vị}
{$400$ đơn vị}
{$4$ đơn vị}
\loigiai{
Chữ số 4 ở hàng chục nghìn có giá trị $40\,000$. Xóa đi thì số đó giảm đúng $40\,000$ đơn vị.
\textbf{Đáp án của câu hỏi là B}
}
\end{ex}

\begin{ex} %[3D1V59]
Số $100\,000$ gồm bao nhiêu chục nghìn?
\choice
{$100$ chục nghìn}
{$1000$ chục nghìn}
{\True $10$ chục nghìn}
{$1$ chục nghìn}
\loigiai{
$100\,000 = 10 \times 10\,000$, tức là 10 chục nghìn.
\textbf{Đáp án của câu hỏi là C}
}
\end{ex}

\begin{ex} %[3D1V59]
Số liền trước của số bé nhất có năm chữ số là
\choice
{\True $9999$}
{$10\,001$}
{$10\,000$}
{$9990$}
\loigiai{
Số bé nhất có năm chữ số là $10\,000$. Số liền trước là: $10\,000 - 1 = 9999$.
\textbf{Đáp án của câu hỏi là A}
}
\end{ex}

\begin{ex} %[3D1V59]
Tìm số có năm chữ số biết chữ số hàng chục nghìn là 7, chữ số hàng trăm gấp đôi chữ số hàng chục nghìn (điều này không thể vì vượt quá 9, ta xét: chữ số hàng nghìn bằng một nửa hàng chục nghìn nếu là số chẵn). Một số có 5 chữ số có chữ số hàng chục nghìn là 8, hàng nghìn bằng nửa hàng chục nghìn, các chữ số còn lại là 0. Số đó là
\choice
{$82\,000$}
{$88\,000$}
{$80\,400$}
{\True $84\,000$}
\loigiai{
Hàng chục nghìn là 8, hàng nghìn là $8 : 2 = 4$, các chữ số còn lại là 0. Số đó là $84\,000$.
\textbf{Đáp án của câu hỏi là D}
}
\end{ex}

\begin{ex} %[3D1V59]
Đọc số $70\,005$ đúng quy tắc là
\choice
{Bảy mươi nghìn năm}
{\True Bảy mươi nghìn không trăm linh năm}
{Bảy vạn năm đơn vị}
{Bảy trăm nghìn linh năm}
\loigiai{
Quy tắc đọc số: hàng trăm bằng 0 đọc là 'không trăm', hàng chục bằng 0 đọc là 'linh': Bảy mươi nghìn không trăm linh năm.
\textbf{Đáp án của câu hỏi là B}
}
\end{ex}

\textbf{Phần 2. Trắc nghiệm đúng sai}
\setcounterpageref{ex}{0}

\begin{ex} %[3D1V59]
Xét tính đúng sai của các nhận định về số có năm chữ số:
\choiceTF
{\True Số 100 000 là số tròn trăm nghìn bé nhất có sáu chữ số.}
{Số lớn nhất có năm chữ số là 100 000.}
{\True Trong số 52 468, chữ số 5 ở hàng chục nghìn có giá trị là 50 000.}
{Số 30 040 đọc là ba mươi nghìn bốn mươi.}
\loigiai{
\begin{itemize}
    \item Ý a) Đúng: Định nghĩa.
    \item Ý b) Sai: Sai vì số lớn nhất có năm chữ số là 99 999.
    \item Ý c) Đúng: Giá trị vị trí hàng.
    \item Ý d) Sai: Sai vì đọc đúng là ba mươi nghìn không trăm bốn mươi.
\end{itemize}
\begin{itemize}
    \item a) \textbf{Đúng}.
    \item b) \textbf{Sai}.
    \item c) \textbf{Đúng}.
    \item d) \textbf{Sai}.
\end{itemize}
}
\end{ex}

\begin{ex} %[3D1V59]
Cho các chữ số $1; 0; 4; 6; 8$:
\choiceTF
{\True Số lớn nhất có 5 chữ số khác nhau ghép từ 5 chữ số trên là 86 410.}
{\True Số bé nhất có 5 chữ số khác nhau ghép từ 5 chữ số trên là 10 468.}
{Chữ số 0 có thể đứng ở hàng chục nghìn để tạo số có 5 chữ số.}
{\True Hiệu giữa số lớn nhất và số bé nhất lập được là: $86\,410 - 10\,468 = 75\,942$.}
\loigiai{
\begin{itemize}
    \item Ý a) Đúng: Thứ tự giảm dần.
    \item Ý b) Đúng: Hàng chục nghìn khác 0, tăng dần.
    \item Ý c) Sai: Sai vì chữ số đầu tiên phải khác 0.
    \item Ý d) Đúng: Phép trừ chính xác.
\end{itemize}
\begin{itemize}
    \item a) \textbf{Đúng}.
    \item b) \textbf{Đúng}.
    \item c) \textbf{Sai}.
    \item d) \textbf{Đúng}.
\end{itemize}
}
\end{ex}

\textbf{Phần 3. Trả lời ngắn}
\setcounterpageref{ex}{0}

\begin{ex} %[3D1V59]
Viết số liền trước của số 60 000.
\shortans{59999}
\loigiai{
\textbf{Đáp số:} $59999$\\
$60\,000 - 1 = 59\,999$.
}
\end{ex}

\begin{ex} %[3D1V59]
Viết số liền sau của số 45 699.
\shortans{45700}
\loigiai{
\textbf{Đáp số:} $45700$\\
$45\,699 + 1 = 45\,700$.
}
\end{ex}

\begin{ex} %[3D1V59]
Tìm giá trị của chữ số 7 trong số 74 521.
\shortans{70000}
\loigiai{
\textbf{Đáp số:} $70000$\\
Chữ số 7 ở hàng chục nghìn có giá trị là 70 000.
}
\end{ex}

\begin{ex} %[3D1V59]
Viết số bé nhất có năm chữ số khác nhau.
\shortans{10234}
\loigiai{
\textbf{Đáp số:} $10234$\\
Số bé nhất có 5 chữ số khác nhau là 10 234.
}
\end{ex}

\begin{ex} %[3D1V59]
Có bao nhiêu số tròn nghìn có năm chữ số từ 20 000 đến 25 000?
\shortans{6}
\loigiai{
\textbf{Đáp số:} $6$\\
Các số: 20 000, 21 000, 22 000, 23 000, 24 000, 25 000 (6 số).
}
\end{ex}

\begin{ex} %[3D1V59]
Tính tổng các chữ số của số bé nhất có 5 chữ số khác nhau.
\shortans{10}
\loigiai{
\textbf{Đáp số:} $10$\\
Số đó là 10 234. Tổng các chữ số: $1 + 0 + 2 + 3 + 4 = 10$.
}
\end{ex}
